\SourceLeaf{58}
\fancyhead[L]{\small REVIEWS OF HIS OWN WORKS}
\pdfbookmark[1]{Theoria motus corporum coelestium. Review of 1809}{anzeige-theoria1809}
\begin{center}\rule{.75\textwidth}{.3pt}\\
Göttingische gelehrte Anzeigen. 17 June 1809.\\
\rule{.75\textwidth}{.3pt}\end{center}
\textit{Theoria motus corporum coelestium in sectionibus conicis Solem ambientium.} By \textsc{Carl Frid. Gauss.} Hamburg, 1809. Published by \textit{Frid. Perthes} and \textit{J. H. Besser.} XII pages of preface, 228 pages of text and 20 pages of tables, with a copperplate. Large quarto.

Many writers are known to have worked on the branch of theoretical astronomy that deals with the parabolic and elliptic motion of celestial bodies. Among them were even mathematicians of the first rank, who devoted works of their own to it. After such predecessors, only a strong reason could justify treating the field anew. Astronomers proper know that reason well enough. It is only for those friends of astronomy who do not make it their main occupation, and who in these last fateful years may have lost sight of the affairs of heaven amid the affairs of earth, that we wish to recall briefly the circumstances that immediately prompted the present work. The problem is to determine with sufficient accuracy the orbit of a celestial body from its motions as observed from Earth over only a moderately long period, knowing nothing more about it than that it moves around the Sun in a conic section according to \textsc{Kepler}’s laws. Until now this problem had never really been
\SourceLeaf{59}tackled in earnest. There are, to be sure, some investigations bearing on it. Without overlooking the ingenuity and analytical skill of which some of them show traces, one must deny the proposed methods everything that practical usefulness requires. At best they give a rough approximation in place of the greatest accuracy possible, or offer a confused heap of undeveloped formulas that would put off even the most tireless calculator in place of a flexible procedure shaped for practical use. Or, finally, they are made wholly useless by errors much smaller than those unavoidable in observation, when they should be applicable to observations such as practical astronomy in its present state allows. This neglect of a problem undoubtedly of great interest in its own right seems to be due in part to the fact that the mathematicians who worked on it were insufficiently familiar with the capabilities and needs of practical astronomy. The main reason, however, is probably that the history of astronomy had yet to present a case where an adequate solution was urgently needed and its benefit would have been keenly felt. Indeed, when \textsc{Kepler}, after discovering his laws, undertook to determine the dimensions of the orbits of the planets known at the time, he had at his disposal both the mean motions, already known with great accuracy, and a wealth of good observations made by \textsc{Tycho} over many years. He had only to select from them whatever he needed for each application of his elegant, though special and comparatively simple, methods. The same resources, indeed greater ones, were available to \textsc{Kepler}’s successors, who, as the art of observation advanced, undertook to determine still more accurately the already fairly well known elements of the planetary orbits.

The case was different with comets, whose motion after \textsc{Newton}’s discoveries was recognized as only a special case of \textsc{Kepler}’s general laws. These celestial bodies are usually visible only briefly. The appearance of their motions depends on the positions that Earth happens to occupy during that time, and the mathematician who sets out to determine the orbits from those appearances
\SourceLeaf{60}has little choice among the observations. He must take them as he finds them, as chance has supplied them. Only very rarely does he have observations made under the circumstances he would wish for in applying special methods. The problem of determining a comet’s parabolic orbit from a few observations was thus as important as it was difficult. \textsc{Newton} himself recognized the difficulties and knew how to overcome them. The zeal and success with which people have pursued this problem from \textsc{Newton}’s time to our own are well known. But there is a fundamental difference between this problem and the one we mentioned above. In the former, the kind of conic section is \textit{prescribed}. In the latter, it must be determined from the infinite variety of possible conic sections, of which the parabola is only one kind, using the observations themselves alone and making no hypothetical assumptions. This restriction understandably simplifies the former problem greatly. For most comets it was a simplification one could and had to allow, since the usually rather low accuracy of the observations and their short duration hardly ever suffice to establish that there is a deviation from a parabola and determine its size. Admittedly, for some comets such a determination was indeed attempted, but, and this is essential, always after a parabolic orbit had already been calculated to serve as an approximation to the ellipse or hyperbola to be determined. Thus here, too, our more general problem remained to some extent dispensable.

Even the discovery of a new enduring celestial body in 1781 did not yet make the need for a solution to this problem felt. Once a parabolic orbit was found inadequate, a circle was tried. Because the planet’s orbit happened to have a rather small eccentricity, the circle represented its motion tolerably well for several years. Its very slow motion, its small distance from the ecliptic and its still fairly bright light meant there was no \textit{danger in delay}. It was found again from one year to the next without any trouble. To determine the deviation from the circle, or the true ellipse, one could therefore wait until the observations could be selected at one’s convenience.

\SourceLeaf{61}Matters were quite different with Ceres, discovered in 1801. This celestial body appears only as a little star of the eighth magnitude. It is difficult to pick it out from the countless host of very similar fixed stars and requires exact knowledge of its position at the time. Its discoverer had observed it for only 41 days. By the time the discovery became known in the rest of Europe, Ceres was already lost in the Sun’s rays and would not be visible again until a year later, in a completely different part of the sky. It was now necessary to predict as accurately as possible where to look for it again. This prediction had to rest solely on the few available observations and rigorous calculation, without uncertain hypotheses. Several astronomers tried the simplest hypothesis, a circular orbit. Admittedly the observations could be brought into only imperfect agreement with it, so whether it would suffice to find the planet again remained at best highly uncertain. Indeed, the outcome later confirmed that by the end of 1801 the circular hypothesis already differed by eleven degrees from the planet’s true position. The astronomers who had the good fortune to find it again themselves declared that recovering it on the basis of so erroneous a hypothesis would have been impossible.

In the summer of 1801, while occupied with something quite different, the author of the present work had come upon some ideas that seemed capable of leading him to a solution of the general problem mentioned above. At another time he might not have immediately pursued and developed these ideas, which initially held only a theoretical appeal for him. But at that very moment \textsc{Piazzi}’s discovery had caught everyone’s attention and his observations had just been made public, so the author could not resist testing the practical applicability of those ideas on them. The outcome of this work is known. It had never been suspected that a short series of observations of a planet could yield a determination of its orbit already very close to the true one and more than accurate enough to find it again after a considerable interval. That possibility was beautifully demonstrated, and the method used proved its worth. Any doubts that might still have remained about how generally useful it was were entirely removed by equally
